\documentclass[12pt,a4paper]{article}
\usepackage[utf8]{inputenc}
\usepackage[T1]{fontenc}
\usepackage[brazil]{babel}
\usepackage{graphicx}
\usepackage{booktabs}
\usepackage{geometry}
\geometry{margin=2cm}
\title{Relatorio de Comparacao das Politicas ABR}
\author{Projeto Final TR2}
\date{24/06/2026}
\begin{document}
\maketitle

\section{Resumo}
Este relatorio compara tres politicas ABR: RATE, BUFFER e HYBRID. A execucao foi realizada com 30 segmentos por politica.

\section{Resultados}
\begin{tabular}{lrrrrrrl}
\toprule
Politica & Throughput & Buffer & Rebuffer & Failover & Trocas & Bitrate & Final\\
\midrule
RATE & 771,52 kbps & 12,58 s & 0 & 0 & 6 & 580,00 kbps & 480p\\
BUFFER & 564,89 kbps & 9,86 s & 1 & 1 & 7 & 590,00 kbps & 720p\\
HYBRID & 1200,33 kbps & 12,88 s & 0 & 0 & 3 & 1066,67 kbps & 1080p\\
\bottomrule
\end{tabular}

\section{Analise}
A politica RATE reage ao throughput, mas oscilou e terminou conservadora em 480p. A politica BUFFER decide apenas pelo nivel do buffer, o que causou rebuffer e failover nesta execucao. A politica HYBRID combinou throughput, buffer e jitter, sustentando 1080p com zero rebuffer.

\section{Failover e mocks}
O failover e implementado pelo ServerManager, que troca para um servidor saudavel quando o download falha. O arquivo mock_failover_policy.py usa MockHealth para validar failover bidirecional e cenarios da politica HYBRID sem depender da rede real.

\end{document}
